The account decides with rule-core's
\begin{Verbatim}[fontsize=\small,xleftmargin=1em]
rule_eval::authorize(rules, records, contexts, choices,
                     signers, ledger, generation)
    -> Result<Vec<SpendRecord>, Refusal>
\end{Verbatim}
over the selected rules, their records, the converted contexts, one local rule
index per context, the claimed signers, the ledger and the grant generation. It
returns the next records or a \code{Refusal}, and neither searches for rules nor
verifies signatures.

\paragraph{The path.} Figure~\ref{fig:fig-check-auth} lists the steps of
\code{\_\_check\_auth}. Four guards run before any rule is read: reserved names,
the classification of a recovery completion, the Protected freeze and a scope
recheck. A recovery completion leaves at step~4 for the recovery controller, so a
recovery completion never depends on the evaluator; enrollment and
reconfiguration are document writes the evaluator decides. Otherwise the account binds
the payload, resolves the selected rule IDs, runs the evaluator, authenticates
every claimed signer and stores each changed record, so a refused batch calls no
verifier. A stale grant refuses with \code{StaleScope} (auth error
\fact{auth-error-stale-scope}) and any other evaluator refusal with
\code{RuleEvalRefused} (\fact{auth-error-rule-eval-refused}), inside
\fact{host-auth-error}; a wallet re-runs the evaluator off chain for the reason.
A deployed verifier that rejects a signature traps in the host's check
(\fact{verifier-trap-error} for Ed25519). Because a controller can be adopted
after a grant is installed, the recheck tests each selected grant rule against
the current excluded set (\code{ScopeEscalation},
\fact{auth-error-scope-escalation}); the install-time property is proved
(\ref{stmt:grant-target-containment}), the recheck tested. Infrastructure
addresses derive from a Wasm hash and registry path fixed at build time, so a
later registry publication cannot change the code an account calls.

\perchfigure{fig-check-auth}

\paragraph{What a signature binds.} With $H$ = SHA-256, the host's payload $P$
for an address credential does not name the account. Perch binds it to the
account's \fact{contract-id-bytes}-byte ID $\mathit{id}_A$ under the tag $\tau$ =
\code{\fact{auth-domain-tag}} and appends the selected rule IDs $\rho$:
\begin{align*}
  P &= H\bigl(\mathrm{XDR}(\mathtt{HashIdPreimage::SorobanAuthorization}\{\mathit{nid}, \mathit{nonce}, \mathit{exp}, \mathit{inv}\})\bigr),\\
  B &= H(\tau \,\|\, \mathit{id}_A \,\|\, P), \qquad
  D = H\bigl(B \,\|\, \mathrm{XDR}(\mathtt{Vec}\langle\mathtt{u32}\rangle\ \rho)\bigr).
\end{align*}
An external signer signs $D$; a delegated signer approves the host's
address-bound payload, which omits $\rho$. OpenZeppelin \fact{oz-v080-tag} signs
$H(\mathrm{XDR}(\mathtt{AuthDigestPreimage}\{A, P, \rho\}))$ for both kinds,
and its \fact{oz-v072-tag} line $H(P \,\|\, \mathrm{XDR}(\rho))$, without the
account (Table~\ref{tab:tab-binding}). A credential's authority is the union of
every live rule that names it; binding the rule IDs stops a relayer, not the
signer, from choosing among those rules. A rule edited in place between signing
and inclusion decides under its new contents, and \code{apply\_grant} takes no
expected version, so two pre-signed installs for one grant ID can land in either
order before they expire.

\perchtable{tab-binding}

\paragraph{Verifiers and hooks.} An external signer pairs key material with a
verifier contract implementing \code{verify(digest, key, signature)}.
Ed25519~\cite{rfc8032} and WebAuthn~\cite{webauthn2} verifiers are deployed; an
\fact{mldsa-verifier}~\cite{fips204} verifier is built and tested. Keys must be in
canonical form, so a threshold, which counts distinct identities, cannot count
one key twice. The WebAuthn verifier does not check the relying-party ID or
origin. An account document names verifiers through the registry; a verified
name denotes provenance its manager admitted, not a proof or an audit. A grant
may name a verifier by address alone; each such verifier may supply at most one
signer per rule, so it cannot cast several threshold votes in one rule, though
it can cast its one vote at will. Before any rule is read, the account refuses
calls on any contract to \fact{reserved-hook-count} reserved names,
\fact{reserved-hook-names}; the first three are OpenZeppelin's policy hooks.
